% AA-38 | Ponowny audyt nocny LVI (2026-11-21); etykiety historyczne zachowane
\chapter{Lemat normalizacyjny Noether}\label{aa:ch:noether-normalization}
\section{Motywacja: algebra skończenie generowana a niezależne parametry}
Niech $A$ będzie skończenie generowaną przemienną algebrą nad ciałem
$K$. Generatory algebry nie muszą być algebraicznie niezależne:
w algebrze ilorazowej $K[X,Y]/(Y^2-X^3)$ klasy $x,y$
zmiennych $X,Y$ są powiązane równaniem $y^2=x^3$. Celem normalizacji Noether jest wykazanie,
że można zawsze wybrać pewną liczbę \emph{niezależnych} parametrów
$y_1,\ldots,y_r\in A$, nad którymi cała algebra jest skończonym
modułem. W dowodzie kluczowa jest zmiana generatorów zmieniająca
wybraną relację na równanie \emph{moniczne}.

Przypomnijmy, że element $a$ algebry $A$ jest \emph{całkowity}
nad podpierścieniem $B\subseteq A$, gdy spełnia równanie
$a^d+c_1a^{d-1}+\cdots+c_d=0$ dla pewnego $d\ge1$
i $c_i\in B$. Jeżeli $A=B[a]$, takie równanie daje
$A=B\cdot1+\cdots+B\cdot a^{d-1}$: wyższe potęgi $a$
redukujemy kolejno, zastępując $a^d$ wyrażeniem
$-c_1a^{d-1}-\cdots-c_d$.

\section{Wybór zmiennych i najwyższej potęgi}
\begin{lemat}[Moniczność po trójkątnej zmianie zmiennych]
\label{aa:lem:noether-monic}
Niech $K$ będzie dowolnym ciałem, $n\ge2$,
a $0\ne f\in K[X_1,\ldots,X_n]$ wielomianem niestałym.
Istnieją dodatnie liczby całkowite $m_1,\ldots,m_{n-1}$
i stała $c\in K\setminus\{0\}$ takie, że po podstawieniu
$X_i=Y_i+Y_n^{m_i}$ ($1\le i<n$), $X_n=Y_n$,
wielomian $c^{-1}f$ staje się moniczny względem $Y_n$,
ze współczynnikami w $K[Y_1,\ldots,Y_{n-1}]$.
Podstawienie jest automorfizmem pierścienia wielomianów.
\end{lemat}
\begin{proof}
\emph{Krok 1: wybór różnych wag.}
Niech $d=\deg f\ge1$ i wybierzmy całkowitą liczbę $M>d+1$.
Dla $1\le i<n$ połóżmy $m_i=M^i$. Każdemu
wielowskaźnikowi $\alpha=(\alpha_1,\ldots,\alpha_n)\in\mathbb N_0^n$
o $|\alpha|\le d$ przypiszmy wagę
$w(\alpha)=\alpha_n+\sum_{i=1}^{n-1}\alpha_i M^i$.
Ponieważ $0\le\alpha_i\le d<M$, liczba $w(\alpha)$
ma w zapisie przy podstawie $M$ cyfry
$\alpha_n,\alpha_1,\ldots,\alpha_{n-1}$.
Jednoznaczność takiego zapisu oznacza, że różne
wielowskaźniki stopnia co najwyżej $d$ mają różne wagi.

\emph{Krok 2: najwyższa potęga dla jednomianu.}
Po wykonaniu podstawienia jednomian
$X^\alpha=\prod_{i=1}^n X_i^{\alpha_i}$
przechodzi w
\[
 Y_n^{\alpha_n}\prod_{i=1}^{n-1}
       (Y_i+Y_n^{m_i})^{\alpha_i}.
\]
Wybranie ze wszystkich nawiasów potęg $Y_n^{m_i\alpha_i}$
daje składnik $Y_n^{w(\alpha)}$ ze współczynnikiem $1$.
Każdy inny składnik zawiera dla pewnego $i<n$
przynajmniej jedną potęgę $Y_i$, a więc ma stopień względem
$Y_n$ ściśle mniejszy od $w(\alpha)$.
Nie ma tu potrzeby dzielenia przez liczby całkowite:
argument działa również wtedy, gdy niektóre współczynniki
dwumianowe znikają w dodatniej charakterystyce.

\emph{Krok 3: brak skasowania najwyższego składnika.}
Zapiszmy $f=\sum_{\alpha}c_\alpha X^\alpha$,
gdzie zbiór indeksów z $c_\alpha\ne0$ jest skończony.
Wybierzmy $\beta$ o największej wadze $w(\beta)$.
Dzięki Krokowi 1 indeks $\beta$ jest jednoznaczny.
Krok 2 pokazuje, że $c_\beta Y_n^{w(\beta)}$ jest
jedynym składnikiem z tak wysoką potęgą $Y_n$ w całym
przekształconym $f$. Jego współczynnik $c_\beta\in K^\times$
nie zależy od $Y_1,\ldots,Y_{n-1}$.
Podzielenie przez $c_\beta$ czyni wielomian monicznym.
Ponieważ $f$ jest niestały, zawiera jednomian o niezerowym
wielowskaźniku, którego waga jest dodatnia. Zatem
$w(\beta)>0$ i uzyskane równanie ma dodatni stopień względem $Y_n$;
bez tego warunku nie wynikałaby całkowitość w zastosowaniu lematu.

\emph{Krok 4: odwracalność podstawienia.}
Przekształcenie $K$-algebry określone przez
$X_i\mapsto Y_i+Y_n^{m_i}$ dla $i<n$, $X_n\mapsto Y_n$
ma odwrotność $Y_i\mapsto X_i-X_n^{m_i}$ i $Y_n\mapsto X_n$.
Zatem jest izomorfizmem $K$-algebr $K[X_1,\ldots,X_n]$
i $K[Y_1,\ldots,Y_n]$, z odwracalną zmianą układu zmiennych.
Nie wprowadza ona nowych relacji ani nie traci istniejących.
\end{proof}

\begin{przyklad}[Dlaczego wybieramy różne wagi]
Niech $f(X_1,X_2)=X_1^2-X_2^3$ i niech $M>4$.
Podstawienie $X_1=Y_1+Y_2^M$, $X_2=Y_2$ daje
\[
 f=Y_1^2+2Y_1Y_2^M+Y_2^{2M}-Y_2^3.
\]
Najwyższa potęga $Y_2$ to $2M$ i jej współczynnik wynosi $1$.
Nawet w charakterystyce $2$, gdy wyraz $2Y_1Y_2^M$
znika, wielomian pozostaje moniczny względem $Y_2$.
Nie wystarcza natomiast dowolna zmiana współrzędnych:
na przykład podstawienie $X_1=Y_1+Y_2$, $X_2=Y_2$
w wielomianie $X_1-X_2$ prowadzi do $Y_1$,
który ma stopień zero względem $Y_2$ i nie daje
monicznego równania dodatniego stopnia dla $Y_2$.
\end{przyklad}

\begin{przyklad}[Jawna normalizacja nad ciałem dwuelementowym]
\label{aa:prz:noether-f2}
Niech $K=\mathbb F_2$ oraz
$A=K[X,Y]/(X^2+XY+Y)$; oznaczmy klasy zmiennych przez
$x,y$. Weźmy $u=x+y^4$. Podstawienie $X=U+V^4$,
$Y=V$ przekształca relację, w charakterystyce $2$, w
\[
 (U+V^4)^2+(U+V^4)V+V
   =V^8+V^5+UV+U^2+V.
\]
Otrzymany wielomian jest moniczny stopnia osiem względem
$V$. W szczególności $y$ spełnia nad $K[u]$ równanie
\[
 y^8+y^5+uy+u^2+y=0,
\]
a więc $A=K[u,y]$ jest generowany jako $K[u]$-moduł
przez $1,y,\ldots,y^7$. Co więcej, $u$ jest
transcendentalny nad $K$: gdyby niezerowy $P\in K[T]$
spełniał $P(u)=0$, to wielomian $P(U)$ należałby
do ideału $(V^8+V^5+UV+U^2+V)$ w $K[U,V]$.
Jest to niemożliwe, ponieważ wielomian tworzący ten ideał
jest moniczny dodatniego stopnia względem $V$,
a każdy jego niezerowy wielokrotny również ma dodatni
stopień względem $V$. Zatem $K[u]\subseteq A$ jest
normalizacją Noether, mimo że ciało bazowe jest skończone.
\end{przyklad}

\section{Twierdzenie o normalizacji}
\begin{twierdzenie}[Lemat normalizacyjny Noether]
\label{aa:tw:noether-normalization}
Niech $K$ będzie ciałem, a $A$ niezerową, przemienną,
skończenie generowaną $K$-algebrą z jedynką,
przy czym odwzorowanie strukturalne $K\to A$ zachowuje jedynkę.
Wówczas istnieją liczba całkowita $r\ge0$ i elementy
$y_1,\ldots,y_r\in A$ takie, że:
\begin{enumerate}[label=(\roman*)]
\item $y_1,\ldots,y_r$ są algebraicznie niezależne nad $K$;
\item $A$ jest skończenie generowanym modułem nad podalgebrą
      $K[y_1,\ldots,y_r]\subseteq A$.
\end{enumerate}
Nie wymagamy, aby $K$ było nieskończone ani aby $A$
była dziedziną całkowitości.
\end{twierdzenie}
\begin{proof}
Niech $n\ge0$ będzie liczbą wybranych generatorów
$A=K[a_1,\ldots,a_n]$. Dowodzimy tezy indukcyjnie względem $n$.

\emph{Krok 1: początek indukcji i przypadek niezależny.}
Jeżeli $n=0$, to $A$ jest obrazem homomorfizmu $K\to A$.
Jądro tego homomorfizmu jest ideałem ciała, a więc jest
albo $(0)$, albo całym $K$. Druga możliwość dawałaby $1_A=0_A$,
co jest sprzeczne z $A\ne0$.
Dlatego $A\cong K$; można przyjąć $r=0$.
Jeśli $n\ge1$ i generatory $a_1,\ldots,a_n$ są algebraicznie
niezależne, to ich podalgebra jest pierścieniem wielomianów,
a $A=K[a_1,\ldots,a_n]$ jest wolnym modułem rzędu jeden
nad tą podalgebrą. Przyjmujemy $r=n$.

\emph{Krok 2: jedna relacja algebraiczna.}
Jeżeli generatory są zależne, to istnieje niezerowy
$f\in K[X_1,\ldots,X_n]$ z $f(a_1,\ldots,a_n)=0$.
Wielomian $f$ nie może być niezerową stałą, ponieważ
$K\hookrightarrow A$ jest injekcją.
Dla $n=1$ dzielimy $f$ przez niezerowy współczynnik
najwyższej potęgi; otrzymujemy moniczne równanie
$a_1^d+c_1a_1^{d-1}+\cdots+c_d=0$ nad $K$.
Zatem $A=K[a_1]$ jest generowany jako $K$-moduł przez
$1,a_1,\ldots,a_1^{d-1}$, co daje tezę z $r=0$.

\emph{Krok 3: redukcja liczby generatorów przy $n\ge2$.}
Zastosujmy Lemat~\ref{aa:lem:noether-monic} do niezerowej
relacji $f$. Dla otrzymanych wykładników $m_i$ definiujemy
$b_i=a_i-a_n^{m_i}$ ($1\le i<n$) i
$B=K[b_1,\ldots,b_{n-1}]\subseteq A$.
Z równości $a_i=b_i+a_n^{m_i}$ wynika
\[
 A=K[a_1,\ldots,a_n]=B[a_n].
\]
Po podstawieniu zmiennych przekształcone równanie $f=0$
przyjmuje, po podzieleniu przez niezerowy skalar,
postać
\[
 a_n^d+c_1(b_1,\ldots,b_{n-1})a_n^{d-1}
 +\cdots+c_d(b_1,\ldots,b_{n-1})=0
\]
z pewnym $d\ge1$, gdzie $c_j\in K[Y_1,\ldots,Y_{n-1}]$.
Zatem $a_n$ jest całkowity nad $B$.
Ten krok nie wymaga, aby $A$ była dziedziną całkowitości:
moniczne równanie pozostaje ważne także w obecności
dzielników zera i elementów nilpotentnych.
Każdy wielomian w $a_n$ można redukować modulo tę relację
do stopnia mniejszego niż $d$.
Dlatego $1,a_n,\ldots,a_n^{d-1}$ generują $A=B[a_n]$
jako $B$-moduł.

\emph{Krok 4: zastosowanie indukcji.}
Algebra $B$ jest niezerową $K$-algebrą generowaną przez
$n-1$ elementów. Z założenia indukcyjnego istnieją
niezależne $y_1,\ldots,y_r\in B$ takie, że $B$ jest
skończonym modułem nad $C=K[y_1,\ldots,y_r]$.
Niech $v_1,\ldots,v_s$ generują $B$ jako $C$-moduł.
Z Kroku 3 wynika, że skończenie wiele elementów
$v_i a_n^j$, gdzie $1\le i\le s$ i $0\le j<d$,
generuje $A$ jako $C$-moduł:
każdy element $A$ jest sumą $B$-wielokrotności $a_n^j$,
a każdy współczynnik z $B$ jest sumą $C$-wielokrotności $v_i$.
Elementy $y_i$ pozostają niezależne w $A$, ponieważ
$B\subseteq A$ jest podpierścieniem i równania
wielomianowe między nimi nie zmieniają się po włączeniu do $A$.
To kończy indukcję i dowód twierdzenia.
\end{proof}

\section{Wymiar transcendencji i przykłady}
\begin{wniosek}\label{aa:wn:noether-dimension}
Jeżeli $A$ jest niezerową, skończenie generowaną dziedziną
całkowitości nad ciałem $K$, to liczba $r$ w normalizacji
Noether jest jednoznacznie wyznaczona i spełnia
\[
 r=\operatorname{trdeg}_K\operatorname{Frac}(A).
\]
\end{wniosek}
\begin{proof}
Niech $B=K[y_1,\ldots,y_r]$ będzie podalgebrą z twierdzenia.
Ponieważ $A$ jest skończonym modułem nad $B$, każdy element
$a\in A$ jest całkowity nad $B$.
Uzasadnijmy to bez dodatkowego założenia: jeżeli
$v_1,\ldots,v_s$ generują $A$ nad $B$, mnożenie przez $a$
wyraża się relacjami
$a v_i=\sum_{j=1}^s c_{ij}v_j$ z $c_{ij}\in B$.
Po zapisaniu $(aI-C)v=0$ i pomnożeniu przez macierz
dołączoną otrzymujemy
$\det(aI-C)v_i=0$ dla każdego $i$.
W szczególności $1=\sum b_i v_i$ dla pewnych $b_i\in B$,
więc $\det(aI-C)\cdot1=0$ w $A$.
Wielomian $\det(TI-C)\in B[T]$ jest moniczny i zeruje się
w $a$. Zatem $a$ jest algebraiczny nad $\operatorname{Frac}(B)$.
Również każdy element $\operatorname{Frac}(A)$ jest ilorazem
dwóch elementów $A$, z których mianownik jest niezerowy;
wszystkie takie ilorazy są algebraiczne nad
$\operatorname{Frac}(B)$, ponieważ elementy algebraiczne
w rozszerzeniu ciał tworzą podciało.
Stąd $\operatorname{Frac}(A)/K(y_1,\ldots,y_r)$ jest
algebraiczne. Niezależność $y_i$ oznacza, że
$\{y_1,\ldots,y_r\}$ jest bazą transcendencji
$\operatorname{Frac}(A)/K$.
Twierdzenie~\ref{aa:tw:transcendence-degree} pokazuje,
że jej liczba elementów nie zależy od wyboru normalizacji.
\end{proof}

\begin{przyklad}[Parabola półsześcienna]
\label{aa:prz:noether-parabola}
Niech $A=K[X,Y]/(Y^2-X^3)$ i niech $x,y$ oznaczają klasy
$X,Y$ w $A$. Relacja $y^2=x^3$ jest moniczna względem $y$.
Dzielenie z resztą w $K[X][Y]$ przez moniczny wielomian
$Y^2-X^3$ daje dla każdego elementu $A$ jednoznaczny zapis
$a(x)+b(x)y$, gdzie $a,b\in K[X]$.
Jednoznaczność wynika stąd, że niezerowy wielomian
$H(X,Y)(Y^2-X^3)$ ma stopień co najmniej $2$ względem $Y$.
Zatem $A=K[x]\oplus K[x]y$ jako $K[x]$-moduł.
Ponadto $x$ jest transcendentalny nad $K$:
gdyby $P(x)=0$ dla niezerowego $P\in K[X]$, to
$P(X)$ należałby do ideału $(Y^2-X^3)$, co jest niemożliwe
z powyższego porównania stopni względem $Y$.
Otrzymujemy normalizację z jednym niezależnym parametrem $x$.
\end{przyklad}

\begin{przyklad}[Hiperbola i nieoczywisty parametr]
Niech $A=K[U,V]/(UV-1)$, a $u,v$ będą klasami zmiennych.
Z relacji $uv=1$ wynika, że $v=u^{-1}$, zatem
$A\cong K[u,u^{-1}]$. Element $u$ jest transcendentalny
nad $K$, bo odwzorowanie $K[U]\to K[U,U^{-1}]$
jest injektywne.
Sam pierścień $K[u]$ nie daje normalizacji: $u^{-1}$
nie jest całkowity nad $K[u]$. Gdyby spełniał moniczne
równanie $(u^{-1})^d+c_1(u^{-1})^{d-1}+\cdots+c_d=0$
ze współczynnikami $c_j\in K[u]$, pomnożenie przez $u^d$
dałoby $1+c_1u+\cdots+c_du^d=0$ w $K[u,u^{-1}]$.
Ponieważ zanurzenie $K[u]\hookrightarrow K[u,u^{-1}]$
jest injektywne, byłaby to tożsamość w $K[u]$.
Podstawienie $u=0$ prowadziłoby do $1=0$ w $K$,
co jest niemożliwe.
Wybierzmy natomiast $s=u+u^{-1}\in A$.
Mamy $u^2-su+1=0$, więc $u$ jest całkowity nad $K[s]$.
Ponadto $u^{-1}=s-u$, stąd $A=K[s,u]$ i
$A=K[s]\cdot1+K[s]\cdot u$.
Element $s$ musi być transcendentalny nad $K$:
gdyby był algebraiczny, to z równania $u^2-su+1=0$
wynikałaby algebraiczność $u$ nad $K$, wbrew powyższemu.
Tak więc $K[s]\subseteq A$ jest normalizacją Noether,
choć $K[u]\subseteq A$ nią nie jest.
\end{przyklad}

\begin{przyklad}[Algebraiczność ciał ułamków nie daje skończoności modułowej]
\label{aa:prz:noether-nonfinite}
Niech $t$ będzie transcendentalny nad $K$.
Rozszerzenie pierścieni $K[t]\subset K(t)$ ma identyczne
ciała ułamków: oba dają $K(t)$, a więc rozszerzenie
ciał ułamków jest algebraiczne stopnia $1$.
Jednak $K(t)$ nie jest skończonym modułem nad $K[t]$.
W przeciwnym razie $1/t$ byłby całkowity nad $K[t]$
(z dowodu Wniosku~\ref{aa:wn:noether-dimension}).
Istniałoby więc równanie
$(1/t)^d+c_1(1/t)^{d-1}+\cdots+c_d=0$ z $c_i\in K[t]$.
Po pomnożeniu przez $t^d$ otrzymujemy
$1+c_1t+\cdots+c_dt^d=0$, co wymusza $1\in tK[t]$;
to niemożliwe, bo każdy wielomian z $tK[t]$ ma wyraz
stały równy zeru.
Kontrprzykład pokazuje, że algebraiczność rozszerzenia
\emph{ciał ułamków} nie wystarcza do skończoności
modułowej rozszerzenia \emph{pierścieni}.
Nie narusza on Twierdzenia~\ref{aa:tw:noether-normalization},
ponieważ $K(t)$ nie jest skończenie generowaną $K$-algebrą.
Aby to udowodnić, przypuśćmy, że
$K(t)=K[f_1(t)/g_1(t),\ldots,f_m(t)/g_m(t)]$,
gdzie $f_i,g_i\in K[t]$, $g_i\ne0$.
Połóżmy $D=g_1\cdots g_m\ne0$.
Wszystkie generatory należą wtedy do $K[t,D^{-1}]$,
więc $K(t)=K[t,D^{-1}]$. Jednak wielomian
$h=1+tD\in K[t]$ jest niestały i względnie pierwszy
z $D$: wspólny dzielnik $h$ i $D$ dzieliłby też
$h-tD=1$. Element $1/h\in K(t)$ nie należy do
$K[t,D^{-1}]$. W przeciwnym razie
$D^N=hP$ dla pewnych $N\ge0$ i $P\in K[t]$;
wówczas $h$ dzieliłby $D^N$, co przeczy względnej
pierwszości $h,D$ i temu, że $h$ nie jest jednostką.
Otrzymana sprzeczność dowodzi zapowiedzianego faktu.
\end{przyklad}



\begin{przyklad}[Normalizacja algebry z elementami nilpotentnymi]
\label{aa:prz:noether-nilpotents}
Niech $A=K[X,E]/(E^2)$, gdzie $K$ jest dowolnym ciałem,
a $x,\varepsilon$ oznaczają klasy zmiennych $X,E$.
Ponieważ $E^2$ jest wielomianem monicznym stopnia dwa
względem $E$, dzielenie w $K[X][E]$ daje jednoznaczny
zapis elementu $A$ jako
$a(x)+b(x)\varepsilon$ dla $a,b\in K[X]$.
W szczególności $A=K[x]\oplus K[x]\varepsilon$ jako
$K[x]$-moduł: ma on bazę $\{1,\varepsilon\}$.
Element $x$ jest transcendentalny nad $K$,
ponieważ niezerowy $P(X)\in K[X]$ nie może być
wielokrotnością $E^2$ w $K[X,E]$.
Zatem $K[x]\subseteq A$ jest normalizacją Noether.
Jednocześnie $\varepsilon^2=0$ i $\varepsilon\ne0$,
więc $A$ nie jest pierścieniem zredukowanym.
Przykład pokazuje, że nie tylko dzielniki zera,
lecz także niezerowe elementy nilpotentne są zgodne
z twierdzeniem o normalizacji.
\end{przyklad}

\begin{przyklad}[Algebra z dzielnikami zera]
\label{aa:prz:noether-zero-divisors}
Niech $A=K[X,Y]/(XY)$, a $x,y$ oznaczają klasy zmiennych.
Element $s=x+y$ jest transcendentalny nad $K$.
Istotnie, homomorfizm $A\to K[X]$ określony przez
$y\mapsto0$, $x\mapsto X$ przesyła $s$ na $X$;
gdyby niezerowy wielomian $P\in K[T]$ zerował się w $s$,
to $P(X)=0$ w $K[X]$, co jest sprzecznością.
Relacja $xy=0$ daje $x(s-x)=0$, czyli
$x^2-sx=0$. Zatem $x$ jest całkowity nad $K[s]$.
Ponieważ $y=s-x$, otrzymujemy $A=K[s,x]$.
Redukcja potęg przez $x^2=sx$ dowodzi
$A=K[s]\cdot1+K[s]\cdot x$, więc $A$ jest skończonym
modułem nad $K[s]$. Jednocześnie $xy=0$, a $x\ne0$
i $y\ne0$: obrazy odpowiednio w ilorazach przez $(y)$
i przez $(x)$ są niezerowymi zmiennymi wielomianowymi.
Przykład potwierdza, że lemat normalizacyjny Noether
nie wymaga braku dzielników zera.
\end{przyklad}
